%Research and development in machine learning heavily rely on deep learning frameworks to provide building blocks and high-level programming interfaces in designing, training and validating a deep neural network. % Please avoid footnotes as much as possible
PyTorch and TensorFlow are the most popular frameworks, backed by Meta AI and Google, respectively.
Both frameworks are similar in many respects, including 1) providing low-level APIs for development, 2) supporting a rich collection of libraries, and 3) maintaining dedicated hubs - PyTorch Hub and TensorFlow Hub - for providing pre-trained ML models. 
% with dedicated hubs, PyTorch Hub and TensorFlow Hub, to provide pre-trained ML models.
%\todo{\textcolor{blue}{have replaced them with\face }Le: please make sure the paper use consistent Hugging Face or HuggingFace. }
\face is a large open-source community that builds tools to enable users to build, train, and deploy machine learning models based on open-source code and technologies.
\face is best known for its \texttt{Transformers} library, which exposes a collection of Python APIs to leverage state-of-the-art deep learning architectures for NLP tasks. 
With the goal to simplify end-to-end NLP tasks, \face \texttt{Transformers} offers a pipeline that performs all pre- and post-processing steps on the given input text data. The overall process of the model inference is encapsulated within these pipelines. With the pipeline, users only need to provide the input texts and the model for the task.  The remaining connections among a model and required pre- and post-processing steps are hidden within the pipeline implementation.

There were various research works and developments to improve the ML ecosystem to be Findable, Accessible, Interoperable, and Reproducible (FAIR).  These existing frameworks aim to make the models, datasets, or both FAIR.  Among these frameworks, HPCFAIR~\cite{verma2021hpcfair} focuses on providing support for model interoperability, search capabilities for datasets and models, and seamless integration into HPC workflows.  The work in \cite{yu2022seamless} extended this work to include support for interoperability across different framework implementations using ONNX and provision to retrain a model with transfer learning.
However, HPCFAIR framework relies on users to handle data pre- and post-processing. In comparison, LM4HPC is equipped to manage data processing within the pipeline design and generate leaderboards for supported HPC tasks. 
%\todo{To Murali: how to differentiate this paper from previous ones? }

%LangChain~\cite{langchain2023} is a framework built to support language models with a goal to provide data-aware connection, between a language model and sources of data, and allow a language model to interact with its environment.
%The core of LangChain is to chain different components to create a effective usage of the LLMs.  LM4HPC APIs proposed in this work can leverage the features from LangChain to enable more sophisticated tasks with LLMs.

General LLMs are trained with data covering general knowledge and information that is usually collected from public domains. 
Domain-specific datasets can be collected for the training of a specialized model or for the fine-tuning of a general-purpose model.
MedQA\cite{jin2020disease} is an example of domain-specific datasets collecting question-answer pairs and textbooks  from professional medical board exams.
ExeBench~\cite{exebench}, another domain-specific dataset for tasks in compilation and software engineering, contains millions of runnable and representative C functions collected from GitHub.
In addition to collecting existing data, ML research has started to automate dataset creation with assistance from the LLMs.  
The developers of LaMini-LM~\cite{wu2023lamini} develop a large set of 2.58M instruction and response pairs based on both existing and newly-generated instructions.  A handful of seed examples from the existing LLM prompts and 2.2M categories from Wikipedia from existing  are submitted to the \texttt{gpt-3.5-turbo} to generate relevant instructions.  Similarly, the responses for the generated instructions are also generated by the \texttt{gpt-3.5-turbo}.

%Mention a few generic and domain-specific QA datasets (e.g. MedQA dataset) and creation methods: 
%https://arxiv.org/abs/2304.14402 TODO

%\todo{Mention a few generic and domain-specific QA datasets (e.g. MedQA dataset) and creation methods: https://arxiv.org/abs/2304.14402}